\section*{Chapitre \ref{ch:g12:diversification-of-life} --- La diversification du vivant}

\begin{solution}{exo:g12:diversification-of-life:1}
La duplication de \omterm{def:g10:universal-dna:gene}{gènes} et leur divergence~; le \omterm{prop:g11:antibiotic-resistance:variation}{transfert horizontal de
gènes}~; l'hybridation avec polyploïdie~; les changements de régulation
des \omterm{prop:g12:diversification-of-life:development}{gènes du développement}~; la \omterm{def:g12:diversification-of-life:symbiosis}{symbiose}~; les comportements transmis.
\end{solution}

\begin{solution}{exo:g12:diversification-of-life:2}
Un ensemble de \omterm{def:g10:universal-dna:gene}{gènes} apparentés descendus d'un même \omterm{def:g10:universal-dna:gene}{gène} ancestral par
duplications successives, chaque copie divergeant ensuite par \omterm{def:g11:mutations:mutation}{mutation}
vers une fonction distincte (\omterm{def:g11:the-eye:photoreceptors}{opsines}, globines).
\end{solution}

\begin{solution}{exo:g12:diversification-of-life:3}
Ses \omterm{def:g10:universal-dna:information}{chromosomes}, un jeu venu de chaque \omterm{def:g10:biodiversity-scales:species}{espèce} parente, n'ont pas
d'homologues avec qui s'apparier à la \omterm{def:g12:meiosis-genetic-shuffling:meiosis}{méiose}, si bien que les gamètes
sont déséquilibrés. Le doublement donne à chaque \omterm{def:g11:cell-cycle-mitosis:chromosome}{chromosome} un
partenaire identique~: l'appariement et la \omterm{def:g12:meiosis-genetic-shuffling:meiosis}{méiose} redeviennent
réguliers et la plante est fertile.
\end{solution}

\begin{solution}{exo:g12:diversification-of-life:4}
Leur propre \omterm{def:g10:universal-dna:information}{ADN} circulaire~; des \omterm{def:g10:cells-common-unit:organelle}{ribosomes} de type bactérien sensibles
aux \omterm{def:g11:antibiotic-resistance:antibiotic}{antibiotiques} antibactériens~; une multiplication par division~;
une double membrane~; des séquences les plus proches d'un groupe de
bactéries actuelles.
\end{solution}

\begin{solution}{exo:g12:diversification-of-life:5}
Un comportement appris des autres membres du groupe et transmis sans
\omterm{def:g10:universal-dna:gene}{gènes}~: le cassage des noix à la pierre chez certaines communautés de
chimpanzés et pas chez d'autres.
\end{solution}

\begin{solution}{exo:g12:diversification-of-life:6}
$14 + 7 = 21$ \omterm{def:g10:universal-dna:information}{chromosomes}~: 7 A venus de l'engrain, 7 A et 7 B venus de
l'amidonnier. Sept paires (A avec A) peuvent se former~; les 7
\omterm{def:g10:universal-dna:information}{chromosomes} B n'ont pas de partenaire et se répartissent au hasard, si
bien que les gamètes sont déséquilibrés~: stérile.
\end{solution}

\begin{solution}{exo:g12:diversification-of-life:7}
La séparation $\alpha$--$\beta$ (50~\% d'identité) est plus ancienne
que la séparation $\beta$--$\gamma$ (80~\%)~: $\alpha$ se détache
d'abord, puis $\beta$ et $\gamma$ se séparent l'un de l'autre.
\end{solution}

\begin{solution}{exo:g12:diversification-of-life:8}
Les \omterm{def:g10:universal-dna:gene}{gènes} des membres sont présents mais le signal qui les allume dans
le flanc de l'embryon n'est pas donné, ou l'est puis s'interrompt~; la
différence est dans la régulation de l'expression, non dans les \omterm{def:g10:universal-dna:gene}{gènes}.
\end{solution}

\begin{solution}{exo:g12:diversification-of-life:9}
Dans la quantité et le calendrier d'un signal de croissance dans le bec
embryonnaire~: une différence de régulation. Un changement de
régulation n'exige qu'une petite \omterm{def:g11:mutations:mutation}{mutation} dans une séquence de
commande, alors qu'une \omterm{def:g11:gene-expression:protein}{protéine} nouvelle en exigerait beaucoup~; il peut
apparaître et être sélectionné en quelques milliers de générations.
\end{solution}

\begin{solution}{exo:g12:diversification-of-life:10}
Sa forme, sa capacité à coloniser la roche nue, sa résistance à la
sécheresse et sa croissance lente n'appartiennent à aucun des deux
partenaires pris seul~; elles émergent de l'association. Le lichen est
un organisme composite doté de propriétés propres.
\end{solution}

\begin{solution}{exo:g12:diversification-of-life:11}
Un comportement transmis~: les matériaux sont disponibles des deux
côtés, les populations sont proches parentes, et la différence suit la
rivière --- une barrière au contact, donc à l'apprentissage, et non aux
\omterm{def:g10:universal-dna:gene}{gènes} ni aux ressources.
\end{solution}

\begin{solution}{exo:g12:diversification-of-life:12}
Un virus a inséré son \omterm{def:g10:universal-dna:gene}{gène} dans le génome d'une \omterm{def:g10:cells-common-unit:cell}{cellule} germinale d'un
ancêtre des primates il y a quelques dizaines de millions d'années~; le
\omterm{def:g10:universal-dna:gene}{gène} a été hérité et conservé. Épreuve~: le même \omterm{def:g10:universal-dna:gene}{gène} devrait occuper la
même position chromosomique chez tous les primates et être absent de
cette position chez les autres mammifères~; sa séquence devrait être la
plus proche de celle du virus.
\end{solution}

\begin{solution}{exo:g12:diversification-of-life:13}
Les \omterm{def:g10:cells-common-unit:organelle}{ribosomes} mitochondriaux sont de type bactérien et fixent les mêmes
\omterm{def:g11:antibiotic-resistance:antibiotic}{antibiotiques}~: à forte dose, le médicament ralentit la synthèse
protéique des \omterm{def:g10:cells-common-unit:organelle}{mitochondries}. Cette sensibilité est héritée de l'ancêtre
bactérien des \omterm{def:g10:cells-common-unit:organelle}{mitochondries}.
\end{solution}

\begin{solution}{exo:g12:diversification-of-life:14}
Le nombre de \omterm{def:g10:universal-dna:information}{chromosomes} du \omterm{prop:g12:diversification-of-life:polyploidy}{polyploïde} ne correspond plus à celui
d'aucun de ses parents, si bien que les croisements avec eux donnent
une descendance stérile~: l'isolement est immédiat. Un nouvel \omterm{def:g10:universal-dna:gene}{allèle}
doit se répandre dans une population au fil de nombreuses générations
avant de distinguer un groupe. Une \omterm{def:g10:biodiversity-scales:species}{espèce} végétale peut donc apparaître
en une seule génération.
\end{solution}

\begin{solution}{exo:g12:diversification-of-life:15}
Au fond, toute séquence nouvelle --- une copie dupliquée qui diverge,
un changement de régulation, les \omterm{def:g10:universal-dna:gene}{gènes} d'un plasmide transféré --- a
commencé quelque part par une \omterm{def:g11:mutations:mutation}{mutation}. Mais la duplication, le
transfert et la polyploïdie déplacent et multiplient des \omterm{def:g10:universal-dna:gene}{gènes} et des
génomes entiers d'un coup~; un changement de régulation produit des
formes nouvelles sans \omterm{def:g10:chemistry-of-life:families}{protéines} nouvelles~; la \omterm{def:g12:diversification-of-life:symbiosis}{symbiose} combine les
génomes d'\omterm{def:g10:biodiversity-scales:species}{espèces} différentes~; et la culture transmet de la variation
sans aucun \omterm{def:g10:universal-dna:gene}{gène}. La \omterm{def:g11:mutations:mutation}{mutation} fournit les lettres~; ces processus
réarrangent des pages et des livres entiers.
\end{solution}

\begin{solution}{pb:g12:diversification-of-life:1}
\textbf{1.} Amidonnier 28, blé tendre 42.

\textbf{2.} Aucune~: les \omterm{def:g10:universal-dna:information}{chromosomes} A et B ne sont pas homologues. Les
gamètes reçoivent des assortiments quelconques des 14, presque jamais
un jeu complet.

\textbf{3.} Avec 28 \omterm{def:g10:universal-dna:information}{chromosomes}, chaque A a un partenaire A identique
et chaque B un B~: 14 paires, \omterm{def:g12:meiosis-genetic-shuffling:meiosis}{méiose} régulière, gamètes équilibrés à 14.

\textbf{4.} Aucun ne peut s'apparier~: 7 \omterm{def:g10:universal-dna:information}{chromosomes} A, 7 B et 7 D,
chacun sans homologue. Stérile.

\textbf{5.} 21 paires~; un grain de pollen porte 21 \omterm{def:g10:universal-dna:information}{chromosomes}.

\textbf{6.} Les deux autres exemplaires fournissent encore la \omterm{def:g11:gene-expression:protein}{protéine}.
Les copies redondantes sont libres d'accumuler des \omterm{def:g11:mutations:mutation}{mutations} et de
diverger --- une duplication à l'échelle d'un génome entier.

\textbf{7.} Dans les champs de la région où l'amidonnier était cultivé,
entre 10\,000 et 8\,000 ans avant nous, l'égilope 2 sauvage poussant en
adventice au milieu de la culture.

\textbf{8.} Que les deux étapes --- hybridation et doublement ---
suffisent à produire le blé tendre à partir de ses parents~: l'histoire
peut être rejouée.

\textbf{9.} Oui~: il ne peut pas donner de descendance fertile avec
l'engrain, donc, selon la définition, c'est une \omterm{def:g10:biodiversity-scales:species}{espèce} distincte, même
s'il en descend.

\textbf{10.} Les \omterm{def:g10:universal-dna:information}{chromosomes} D de l'égilope 2 sont homologues du jeu D
du blé et s'apparient avec eux, si bien qu'un croisement suivi de
rétrocroisements peut y transférer le \omterm{def:g10:universal-dna:gene}{gène}~; les \omterm{def:g10:universal-dna:information}{chromosomes} d'une
plante non apparentée ne s'apparieraient pas du tout.

\textbf{11.} Les \omterm{def:g10:universal-dna:gene}{gènes} eux-mêmes ne diffèrent presque pas~; la
différence entre les corps doit donc tenir à la façon dont les \omterm{def:g10:universal-dna:gene}{gènes}
sont utilisés.

\textbf{12.} Un changement de régulation de l'expression --- l'étendue
de la région, le long de l'axe, où le \omterm{def:g10:universal-dna:gene}{gène} est actif.

\textbf{13.} Les \omterm{def:g10:universal-dna:gene}{gènes} des membres sont présents et commencent à agir~;
le signal qui entretient la croissance du membre manque, et les
bourgeons s'arrêtent. Là encore, la régulation, et non la perte des
\omterm{def:g10:universal-dna:gene}{gènes}.

\textbf{14.} Elle montre que changer le lieu d'expression du \omterm{def:g10:universal-dna:gene}{gène}
suffit à changer la région du corps~: la correspondance entre
expression et forme est causale.

\textbf{15.} Le changement du blé tendre a ajouté des jeux entiers au
génome (sa taille)~; celui du python a changé le lieu d'utilisation de
\omterm{def:g10:universal-dna:gene}{gènes} communs (leur régulation).

\textbf{16.} \omterm{def:g10:universal-dna:information}{ADN} circulaire propre (attendu~: aucun, tous les \omterm{def:g10:universal-dna:gene}{gènes}
dans le \omterm{def:g10:cells-common-unit:organelle}{noyau})~; \omterm{def:g10:cells-common-unit:organelle}{ribosomes} bactériens (attendu~: des \omterm{def:g10:cells-common-unit:organelle}{ribosomes}
eucaryotes)~; multiplication par division (attendu~: un assemblage à
partir de pièces)~; double membrane (attendu~: une membrane unique
comme les autres \omterm{def:g10:cells-common-unit:organelle}{organites}). Et leurs \omterm{def:g10:universal-dna:gene}{gènes} ressemblent à des \omterm{def:g10:universal-dna:gene}{gènes}
bactériens (attendu~: des \omterm{def:g10:universal-dna:gene}{gènes} de type nucléaire).

\textbf{17.} Dans le \omterm{def:g10:cells-common-unit:organelle}{noyau}~: avec le temps, la plupart des \omterm{def:g10:universal-dna:gene}{gènes} de la
bactérie ont été transférés au génome nucléaire, et le génome propre de
l'\omterm{def:g10:cells-common-unit:organelle}{organite} s'est réduit à un vestige.

\textbf{18.} Les \omterm{def:g10:universal-dna:gene}{gènes} mitochondriaux ne passent que de la mère à
l'enfant~; l'\omterm{def:g10:universal-dna:information}{ADN} mitochondrial d'une personne retrace la lignée
maternelle sans interruption.

\textbf{19.} La \omterm{def:g10:cells-common-unit:organelle}{mitochondrie} a cédé la plupart de ses \omterm{def:g10:universal-dna:gene}{gènes} à son hôte
et ne peut plus exister hors de la \omterm{def:g10:cells-common-unit:cell}{cellule}~; les partenaires du lichen
gardent leurs génomes et peuvent, difficilement, être séparés.

\textbf{20.} L'engrain (14, A), l'égilope 1 (14, B), l'égilope 2 (14,
D)~: 42 \omterm{def:g10:universal-dna:information}{chromosomes}~; la «~régulation~»~; la \omterm{def:g10:cells-common-unit:organelle}{mitochondrie}.
\end{solution}
